% "Program summary" block required by CPC (Computer Programs in Physics)
\noindent\begin{tabularx}{\linewidth}{@{}>{\bfseries}p{0.30\linewidth}X@{}}
Program title: & \BMR\ (Brillouin-zone Utilities for Materials Research) \\
CPC Library link to program files: & (to be assigned by the journal) \\
Developer's repository link: & \url{https://github.com/brunohdmacedo/brumare} \\
Web application: & \url{https://brunohdmacedo.tech/brumare/} \\
Licensing provisions: & GNU General Public License v3.0 or later \\
Programming language: & JavaScript (ECMAScript 2017, browser), Python 3.8+ \\
Supplementary material: & tutorial in six languages (HTML and PDF), example inputs and outputs \\
Nature of problem: & Analysis of electronic-structure outputs (eigenvalues, orbital projections, structures, energies, charge densities) produced by heterogeneous codes and file formats, without recomputing the Fermi energy, with consistent units and with publication-ready, citable outputs. \\
Solution method: & Code-specific readers convert every output to a common internal representation (energies in Hartree, reciprocal lattice in Bohr$^{-1}$, structures in \AA); a single analysis layer implements Gaussian-broadened densities of states, projection grouping, electron-counting band edges, parabolic effective masses, minimum-image structural matching, polyhedral distortion indices, defect thermodynamics, trilinear density interpolation, Voronoi and on-grid Bader partitions and density-difference descriptors. \\
Additional comments: & Runs in any modern browser (no data leave the user's computer), as a desktop application (pywebview, Microsoft Edge WebView2 on Windows) and on clusters (numpy, matplotlib). Interface, help and tutorial available in Portuguese, English, Spanish, French, Italian and Latin. \\
\end{tabularx}
